\section{Schema Matching}
\label{sec:schema_matching}

Schema matching identifies correspondences between attributes in source datasets and the target schema~\cite{rahm2001survey, zhang2023schema}. This section describes how schema matching is handled in the LLM-based and the human pipeline and compares the results of both approaches.

\textbf{Human Pipeline.} As the source schemata are relatively small (5--10 attributes per dataset), the data engineers manually inspected column names and sample values to determine mappings to the target schema. This produced 67 attribute correspondences across all three case studies. For comparison, we also run two non-LLM schema matching methods provided by PyDI: A label-based matcher which compares attribute labels using Monge-Elkan similarity. The best-performing internal similarity metric (Levenshtein or Jaro-Winkler) is automatically selected per use case. An instance-based matcher which tokenizes column values by whitespace, computes TF-IDF vectors, and compares source and target column vectors using cosine similarity. We use the fusion validation set (well-known entities with ground truth values, see Section~\ref{sec:data_fusion}) as the target column representation.

\textbf{LLM Pipeline.} \label{sec:sm_approach}
The LLM pipeline directly prompts GPT-5.2 to perform schema matching. Each source dataset is matched to the target schema using a single LLM call. The prompt consists of a system message instructing the model to act as a schema alignment expert, followed by the source table representation and the target schema. The source representation includes column names, a value summary (unique counts and example values per column), and sample rows rendered as a markdown table. Rows are selected by completeness to maximize informative content. The target schema is provided as a JSON Schema document specifying attribute names, data types, and attribute descriptions. GPT-5.2 returns a JSON array of correspondences, pairing each source column with its target column or indicating no match. For cases where column names are non-descriptive (e.g., ``Attribute\_1'' in the FullContact and MusicBrainz datasets), the prompt explicitly instructs the model to infer attribute semantics from the data values.
The complete prompt is found in the accompanying repository.\footnote{\url{https://github.com/wbsg-uni-mannheim/automatic-data-integration/tree/main/prompts}\label{fn:prompts}}

% Schema matching results comparing Human, LLM, and heuristic matchers

\begin{table}[t]
\caption{Schema matching F1 scores (macro-averaged across datasets).}
\label{tab:sm_results}
\centering
\footnotesize
\begin{tabular}{@{}lcccc@{}}
\toprule
\textbf{Use Case} & \textbf{Human} & \textbf{LLM} & \textbf{Label-based} & \textbf{Instance-based} \\
\midrule
Games & 1.00 & 1.00 & .23 & .00 \\
Companies & 1.00 & 1.00 & .35 & .32 \\
Music & 1.00 & 1.00 & .38 & .24 \\
\midrule
\textit{Average} & 1.00 & 1.00 & .32 & .19 \\
\bottomrule
\end{tabular}
\end{table}


\textbf{Results.} \label{sec:sm_results}
Table~\ref{tab:sm_results} reports schema matching F1 scores macro-averaged across source datasets for all four approaches.

Both the human data engineers and GPT-5.2 correctly identified all 67 attribute correspondences, achieving perfect F1 scores. For the LLM, this includes two datasets with non-descriptive headers (``Attribute\_1'', ``Attribute\_2'', etc.) where the model differentiated columns based solely on value patterns, for instance distinguishing assets from revenue by value magnitudes. The label-based matcher achieved high precision but low recall (average F1=0.32), correctly matching obvious name pairs (e.g., ``tracks\_track-name'' to ``tracks'') but missing semantically equivalent attributes with different names (e.g., ``imprint'' to ``label''). The instance-based matcher performs worse (average F1=0.19) and fails completely on Games (F1=0.00) as data values do not overlap sufficiently between source and target.

 
\textbf{Discussion.}\label{sec:sm_limitations} GPT-5.2 was able to discover the same schema correspondences that were manually determined by the data engineers. A major limitation of the applied matching technique is that it assumes single tables and does not cover scenarios where entity data is distributed over several normalized tables. It also does not cover more complex correspondences (1:n or higher-order) as well as data aggregation. If required, approaches to further improve schema matching could include applying duplicate-based schema matching~\cite{naumann2006dataFusionThreeSteps}, as LLMs are quite capable of identifying entity correspondences between datasets (see Section~\ref{sec:em_approach}), or using more advanced LLM-based schema matching techniques~\cite{liu2025magneto}.

